\chapter{Huấn luyện và đánh giá}
\label{chap:results}

Chương này trình bày quy trình huấn luyện thực tế, kết quả định lượng trên cả hai bộ dữ liệu, các
nghiên cứu loại trừ (ablation), phân tích định tính và phần thảo luận. Mọi con số đều được lấy trực
tiếp từ kết quả thực nghiệm trên hai bộ dữ liệu.

% \section{Quy trình huấn luyện theo bốn giai đoạn}

% Bốn giai đoạn được thực hiện liên hoàn, dùng chung một cấu hình thống nhất, theo thứ tự phụ thuộc
% minh hoạ ở Hình~\ref{fig:phase-flow}.

% \begin{figure}[H]
% \centering
% \begin{tikzpicture}[
%   font=\small,
%   ph/.style={rectangle, rounded corners, draw, thick, align=center, minimum height=1.0cm, text width=3.0cm},
%   arr/.style={-{Stealth[length=2.5mm]}, thick},
% ]
%   \node[ph, fill=blue!8] (p1) {\textbf{Phase 1}\\Train từ đầu\\\scriptsize sinh mô hình CMO};
%   \node[ph, fill=orange!12, below left=1.0cm and -0.5cm of p1] (p0) {\textbf{Phase 0}\\Tái dùng checkpoint};
%   \node[ph, fill=green!12, below right=1.0cm and -0.5cm of p1] (p3) {\textbf{Phase 3}\\Nghiên cứu nguồn\\tri thức (cần CMO)};
%   \node[ph, fill=green!8, right=1.4cm of p1] (p2) {\textbf{Phase 2}\\Thích nghi CLIP\\\scriptsize (chỉ cần dataset)};
%   \draw[arr] (p1) -- (p0);
%   \draw[arr] (p1) -- (p3);
%   \node[above=0.2cm of p2, font=\scriptsize, text=green!40!black] {song song};
% \end{tikzpicture}
% \caption{Phụ thuộc giữa các giai đoạn. Phase 1 sinh mô hình CMO làm đối chứng cho Phase 0 và Phase
% 3; Phase 2 độc lập (chỉ cần dataset) nên chạy song song được.}
% \label{fig:phase-flow}
% \end{figure}

% Mỗi lần chạy đều lưu lại cấu hình, lịch sử huấn luyện theo từng epoch, mô hình tốt nhất cùng các
% hình minh hoạ (confusion matrix, t-SNE); cuối cùng, toàn bộ kết quả được tổng hợp thành bảng so sánh
% sắp theo balanced accuracy.

\section{Kết quả định lượng}

Để tiện theo dõi, kết quả được trình bày \emph{theo từng track}, và \textbf{mỗi bảng đối chiếu trực
tiếp hai bộ dữ liệu} CIFAR-100-LT và CUB-200-LT cạnh nhau. Với mỗi phương pháp, nhóm báo cáo hai
thước đo chính: balanced accuracy (bal-acc) và độ chính xác trên nhóm few-shot (đuôi).

\subsection{Track 1 --- Huấn luyện từ đầu}

Bảng~\ref{tab:scratch} đối chiếu nhóm huấn luyện-từ-đầu trên cả hai bộ dữ liệu. Trên CIFAR, bức
tranh rất rõ theo tầng can thiệp: Baseline chỉ đạt few-shot $0.079$ (đuôi gần như chết), can thiệp ở
\emph{mất mát} (Balanced Softmax) nâng lên $0.181$, và can thiệp đồng thời ở \emph{dữ liệu} và
\emph{mất mát} (CMO) đạt cao nhất ($0.468$ bal-acc, $0.299$ few-shot --- gần 4 lần Baseline). Trần
của việc học-từ-đầu trên CIFAR vì thế chỉ vào khoảng $0.47$ bal-acc.

Trên CUB (fine-grained) thì mọi phương pháp from-scratch đều \emph{sụp đổ} về mức $0.18$--$0.28$
bal-acc, và đáng chú ý là \textbf{CMO không còn dẫn đầu} ($0.195$, thậm chí dưới cả Baseline
$0.258$): cơ chế dán-mảnh thô của CMO không phù hợp với bài toán phân biệt loài chim ở những chi
tiết tinh vi. Dù xếp hạng nội bộ có khác nhau, kết luận chung là như nhau --- học-từ-đầu chạm trần
rất thấp, đặc biệt trên miền khó, nên việc \emph{mượn tri thức ngoài} là cần thiết.

\begin{table}[H]
\centering
\caption{Track 1 --- huấn luyện từ đầu, đối chiếu hai bộ dữ liệu. Sắp theo bal-acc trên CIFAR.}
\label{tab:scratch}
\begin{tabular}{lcccc}
\toprule
 & \multicolumn{2}{c}{\textbf{CIFAR-100-LT}} & \multicolumn{2}{c}{\textbf{CUB-200-LT}} \\
\cmidrule(lr){2-3}\cmidrule(lr){4-5}
\textbf{Phương pháp} & bal-acc & few & bal-acc & few \\
\midrule
\textbf{CMO} & \textbf{0.468} & \textbf{0.299} & 0.195 & 0.120 \\
Balanced Softmax & 0.441 & 0.181 & 0.271 & 0.197 \\
Decoupling (cRT) & 0.407 & 0.094 & 0.261 & 0.169 \\
Baseline & 0.401 & 0.079 & 0.258 & 0.160 \\
SupCon & 0.367 & 0.021 & 0.183 & 0.099 \\
Meta (ProtoNet) & 0.351 & 0.064 & 0.276 & 0.198 \\
\bottomrule
\end{tabular}
\end{table}

SupCon kém kỳ vọng ở đuôi (CIFAR few $0.021$) vì với 5 ảnh/lớp, số cặp dương quá ít để hình thành
cụm tương phản tốt; Meta (ProtoNet) yếu ở $100$-way như đã dự đoán. Cần lưu ý: nhóm from-scratch
trên CUB được chấm khi ngưỡng nhóm-shot còn ở mặc định, nên \emph{cách chia many/medium/few không
hợp lệ} (xem phần thảo luận); balanced accuracy không phụ thuộc cách chia nhóm nên vẫn dùng để so
sánh được, còn cột few của CUB chỉ mang tính tham khảo.

\subsection{Track 2 --- Tái dùng checkpoint}

Bảng~\ref{tab:reuse} đối chiếu các kỹ thuật tái dùng (chạy trên chính các checkpoint từ-đầu, nên
nhất quán với Track 1). Trên \emph{cả hai} bộ, VLM fusion --- trộn mô hình thị giác với CLIP
zero-shot --- dẫn đầu ($0.671$ CIFAR / $0.483$ CUB) và vượt từng chuyên gia thành phần, cho thấy giá
trị của tri thức ngôn ngữ ngay cả khi mới chỉ ``trộn''. Các kỹ thuật thuần-thị-giác (ensemble,
tier-fusion, $\tau$-norm) chỉ cải thiện nhẹ và chủ yếu ở head. Riêng CLIP zero-shot (không huấn
luyện gì) đã đạt $\approx 0.63$ / $0.46$, vượt mọi mô hình thị giác từ-đầu --- một tín hiệu sớm cho
hướng tri thức ngoài ở Track 3.

\begin{table}[H]
\centering
\caption{Track 2 --- tái dùng checkpoint (trên checkpoint từ-đầu), đối chiếu hai bộ dữ liệu.}
\label{tab:reuse}
\begin{tabular}{lcccc}
\toprule
 & \multicolumn{2}{c}{\textbf{CIFAR-100-LT}} & \multicolumn{2}{c}{\textbf{CUB-200-LT}} \\
\cmidrule(lr){2-3}\cmidrule(lr){4-5}
\textbf{Kỹ thuật} & bal-acc & few & bal-acc & few \\
\midrule
\textbf{VLM fusion} (CLIP $+$ vision) & \textbf{0.671} & 0.614 & \textbf{0.483} & 0.444 \\
CLIP zero-shot & 0.627 & 0.638 & 0.460 & 0.445 \\
$\tau$-norm & 0.449 & 0.313 & 0.195 & 0.162 \\
Ensemble $+$ TTA & 0.448 & 0.149 & 0.289 & 0.198 \\
Tier-fusion & 0.425 & 0.132 & 0.200 & 0.114 \\
\bottomrule
\end{tabular}
\end{table}

Vì bảng này dùng đúng các checkpoint từ-đầu (không phải bản tiền huấn luyện), các con số nằm nhất
quán giữa Track 1 và Track 3: ví dụ CMO ở đây cho $\approx 0.46$ bal-acc trên CIFAR, khớp với cả
leaderboard từ-đầu lẫn mốc đối chứng ở Track 3.

\subsection{Track 3 --- Thích nghi mô hình nền tảng (kết quả chính)}

Đây là phần quan trọng nhất (Bảng~\ref{tab:foundation}). Track tri thức ngoài tạo ra một bước nhảy
quyết định trên \emph{cả hai} bộ: từ trần from-scratch $0.47$ lên $\mathbf{0.82}$ bal-acc trên
CIFAR, và từ $0.20$--$0.28$ lên $\mathbf{0.85}$ trên CUB --- khoảng cách càng lớn ở miền khó hơn
(CUB). Trên cả hai, mô hình nền tảng thị giác thứ hai (DINOv2-LIFT) là expert đơn mạnh nhất
($0.811$ / $0.841$) và vượt rõ LIFT-CLIP ($0.718$ / $0.670$); fusion gộp các expert đạt cao nhất
($0.823$ / $0.849$) mà không tụt ở bất kỳ nhóm-shot nào. DINOv2-LIFT thậm chí còn mạnh hơn trên CUB,
phù hợp với việc DINOv2 (patch nhỏ $14$, tự giám sát) bắt chi tiết fine-grained rất tốt.

\begin{table}[H]
\centering
\caption{Track 3 --- thích nghi mô hình nền tảng, đối chiếu hai bộ dữ liệu. Sắp theo bal-acc trên CIFAR.}
\label{tab:foundation}
\begin{tabular}{llcccc}
\toprule
 & & \multicolumn{2}{c}{\textbf{CIFAR-100-LT}} & \multicolumn{2}{c}{\textbf{CUB-200-LT}} \\
\cmidrule(lr){3-4}\cmidrule(lr){5-6}
\textbf{Expert} & \textbf{Nguồn} & bal-acc & few & bal-acc & few \\
\midrule
\textbf{Fusion (greedy)} & gộp & \textbf{0.823} & \textbf{0.755} & \textbf{0.849} & \textbf{0.864} \\
Fusion (theo nhóm-shot) & gộp & 0.820 & 0.753 & 0.846 & 0.854 \\
\textbf{DINOv2-LIFT} & thị giác thứ hai & 0.811 & 0.737 & 0.841 & 0.859 \\
LIFT (CLIP) & vision-language & 0.718 & 0.635 & 0.670 & 0.593 \\
Tip-Adapter & truy hồi & 0.668 & 0.678 & 0.565 & 0.516 \\
Tip-Adapter-F & truy hồi (FT) & 0.643 & 0.400 & 0.649 & 0.529 \\
CLIP-LLM & ngôn ngữ & 0.626 & 0.658 & 0.484 & 0.466 \\
CLIP zero-shot & --- & 0.627 & 0.638 & 0.460 & 0.445 \\
CMO (đối chứng) & nội tại & 0.462 & 0.307 & 0.195 & 0.120 \\
\bottomrule
\end{tabular}
\end{table}

% Hình~\ref{fig:fewshot-progression} trực quan hoá thông điệp cốt lõi trên CIFAR: few-shot accuracy
% \emph{tăng đơn điệu} qua các mốc đại diện cho mạch tiếp cận.

% \begin{figure}[H]
% \centering
% \begin{tikzpicture}
% \begin{axis}[
%   ybar, width=0.85\textwidth, height=6cm,
%   bar width=20pt,
%   ymin=0, ymax=0.85,
%   ylabel={Few-shot accuracy (đuôi)},
%   symbolic x coords={Baseline,CMO,CLIP-zeroshot,LIFT-CLIP,DINOv2-LIFT,Fusion},
%   xtick=data, x tick label style={rotate=20, anchor=east, font=\small},
%   nodes near coords, nodes near coords style={font=\scriptsize},
%   enlarge x limits=0.12,
% ]
% \addplot[fill=blue!40] coordinates {
%   (Baseline,0.079) (CMO,0.299) (CLIP-zeroshot,0.638) (LIFT-CLIP,0.635) (DINOv2-LIFT,0.737) (Fusion,0.755)
% };
% \end{axis}
% \end{tikzpicture}
% \caption{CIFAR-100-LT: few-shot accuracy tăng đơn điệu từ Baseline ($0.079$) $\to$ CMO (trần
% from-scratch, $0.299$) $\to$ track tri thức ngoài (CLIP $0.638$, LIFT, DINOv2 $0.737$, Fusion
% $0.755$). Đuôi được cứu chủ yếu nhờ \emph{mượn tri thức ngoài}.}
% \label{fig:fewshot-progression}
% \end{figure}

\section{Nghiên cứu loại trừ (Ablation Study)}

\subsection{Ablation A --- đóng góp của pipeline so với backbone (DINOv2)}

Câu hỏi đặt ra là: trong kết quả cuối cùng, bao nhiêu phần đến từ \emph{backbone mạnh} (DINOv2) và
bao nhiêu đến từ \emph{pipeline của nhóm} (loss nhận-biết-đuôi $+$ adapter/cosine head)? Để tách
bạch, nhóm dựng một chuỗi ba mốc trên \emph{cùng} đặc trưng DINOv2
(Bảng~\ref{tab:ablationA}). Trên CIFAR, chỉ riêng việc đổi cross-entropy sang Balanced Softmax đã
kéo few-shot từ $0.020$ lên $0.325$ (đóng góp của \emph{loss}); thêm adapter và cosine init đẩy tiếp
lên $0.737$ (đóng góp của \emph{kiến trúc}). Trên CUB hiệu ứng còn rõ rệt hơn
($0.077\to0.675\to0.859$). Kết luận rút ra là \emph{backbone mạnh thôi chưa đủ}: nếu chỉ gắn một đầu
tuyến tính ngây thơ thì đuôi vẫn sụp, và phần lớn giá trị thực sự đến từ pipeline tail-aware của
nhóm.

\begin{table}[H]
\centering
\caption{Ablation A trên đặc trưng DINOv2: tách đóng góp của loss và kiến trúc.}
\label{tab:ablationA}
\begin{tabular}{lcccc}
\toprule
& \multicolumn{2}{c}{\textbf{CIFAR-100-LT}} & \multicolumn{2}{c}{\textbf{CUB-200-LT}} \\
\textbf{Mốc} & bal-acc & few & bal-acc & few \\
\midrule
DINOv2 $+$ Linear (CE) & 0.463 & 0.020 & 0.490 & 0.077 \\
DINOv2 $+$ Linear (Balanced-Softmax) & 0.657 & 0.325 & 0.744 & 0.675 \\
\textbf{DINOv2-LIFT} ($+$adapter $+$cosine init) & \textbf{0.811} & \textbf{0.737} & \textbf{0.841} & \textbf{0.859} \\
\bottomrule
\end{tabular}
\end{table}

\subsection{Ablation B --- độ sâu fine-tune CLIP: fine-tune nặng làm hại đuôi}

Đây là ablation nhằm biện minh cho thiết kế \emph{đóng băng}. Nhóm fine-tune CLIP ở ba độ sâu tăng
dần --- chỉ head, rồi block cuối, rồi toàn bộ visual encoder --- với Balanced Softmax, và đối chiếu
trên cả hai bộ dữ liệu (Bảng~\ref{tab:ablationB}). Trên CIFAR, balanced accuracy có tăng theo độ sâu
($0.461\to0.536\to0.558$) nhưng few-shot lại tệ đi và luôn ở mức rất thấp (cao nhất chỉ $0.128$);
G-mean của các cấu hình này chỉ cỡ $0.004$--$0.007$, cho thấy rất nhiều lớp tail bị ``chết''. Trên
CUB, hiện tượng còn rõ hơn: fine-tune \emph{toàn bộ} encoder thậm chí cho bal-acc thấp nhất
($0.411$, dưới cả linear-probe $0.438$). So với LIFT \emph{đóng băng} (few-shot $0.635$ / $0.593$),
mọi cấu hình fine-tune đều thua xa ở đuôi trên cả hai bộ. Kết quả xác nhận đúng luận điểm của LIFT:
với 5 ảnh/lớp, lan truyền gradient vào backbone gây overfit và xói mòn tri thức zero-shot ở các lớp
hiếm.

\begin{table}[H]
\centering
\caption{Ablation B --- độ sâu fine-tune CLIP (Balanced Softmax), đối chiếu hai bộ dữ liệu. Fine-tune
càng sâu, đuôi càng tệ; LIFT đóng băng vượt trội.}
\label{tab:ablationB}
\begin{tabular}{lcccc}
\toprule
 & \multicolumn{2}{c}{\textbf{CIFAR-100-LT}} & \multicolumn{2}{c}{\textbf{CUB-200-LT}} \\
\cmidrule(lr){2-3}\cmidrule(lr){4-5}
\textbf{Cấu hình} & bal-acc & few & bal-acc & few \\
\midrule
Chỉ head (linear probe) & 0.461 & 0.086 & 0.438 & 0.146 \\
Head $+$ block cuối & 0.536 & 0.128 & 0.548 & 0.266 \\
Toàn bộ visual encoder & 0.558 & 0.111 & 0.411 & 0.084 \\
\midrule
\textbf{LIFT (đóng băng)} & \textbf{0.718} & \textbf{0.635} & \textbf{0.670} & \textbf{0.593} \\
\bottomrule
\end{tabular}
\end{table}

\subsection{Augment đặc trưng: sinh dữ liệu và mixup không giúp}

Khi so sánh ba biến thể trên cùng nền LIFT-CLIP (Bảng~\ref{tab:augment}), cả \emph{sinh đặc trưng
bằng diffusion} lẫn \emph{feature mixup} đều \emph{thấp hơn} LIFT-CLIP gốc trên \emph{cả hai}
dataset. Đây là một kết quả âm trung thực: khi đặc
trưng nền tảng đã đủ mạnh, việc augment thêm chủ yếu đưa vào nhiễu thay vì thông tin mới. Riêng
diffusion còn làm tail tệ đi rõ rệt trên CIFAR (few $0.505$).

\begin{table}[H]
\centering
\caption{Augment đặc trưng so với LIFT-CLIP gốc (balanced accuracy).}
\label{tab:augment}
\begin{tabular}{lcc}
\toprule
\textbf{Phương pháp} & \textbf{CIFAR-100-LT} & \textbf{CUB-200-LT} \\
\midrule
LIFT-CLIP (gốc) & \textbf{0.717} & \textbf{0.670} \\
LIFT-CLIP $+$ feature mixup & 0.698 & 0.654 \\
LIFT-CLIP $+$ diffusion & 0.675 & 0.655 \\
\bottomrule
\end{tabular}
\end{table}

% \section{Kết quả định tính}

% Một ``slide demo đuôi'' điển hình minh hoạ rất rõ giá trị của tri thức ngoài: với một ảnh thuộc lớp
% hiếm, mô hình thị giác from-scratch thường dự đoán sai (gán nhầm sang một lớp head trông tương tự),
% trong khi LIFT/DINOv2 --- nhờ khởi tạo ngữ nghĩa và loss nhận-biết-đuôi --- lại dự đoán đúng. Các
% trường hợp như vậy cho thấy bằng trực giác vì sao tri thức ngoài cứu được đuôi. Những hình tương ứng
% (confusion matrix, t-SNE của đặc trưng) được sinh ra từ pipeline đánh giá và có thể chèn vào các vị
% trí dưới đây.

% \begin{figure}[H]
% \centering
% \fbox{\parbox[c][3.2cm][c]{0.8\textwidth}{\centering [Hình minh hoạ: kết quả định tính]\\[2pt]
% \small Ví dụ dự đoán trên lớp tail: vision (sai) vs CLIP/LIFT/DINOv2 (đúng).}}
% \caption{Phân tích định tính trên các lớp đuôi (placeholder --- chèn ảnh từ pipeline đánh giá).}
% \label{fig:qualitative}
% \end{figure}

% \begin{figure}[H]
% \centering
% \fbox{\parbox[c][3.2cm][c]{0.8\textwidth}{\centering [Hình minh hoạ: so sánh confusion matrix]\\[2pt]
% \small Confusion matrix chuẩn hoá: Baseline (tail ``chết'') vs DINOv2-LIFT (tail sống).}}
% \caption{So sánh confusion matrix: phương pháp từ đầu bỏ rơi đuôi, track nền tảng phân bố đều hơn.}
% \label{fig:confmat}
% \end{figure}

\section{Thảo luận}

\subsection{Các phát hiện chính (nhất quán trên cả hai dataset)}
\begin{enumerate}
  \item \textbf{Tri thức ngoài là bước nhảy quyết định.} Trên CIFAR, kết quả đi từ from-scratch
        $0.47$ lên nền tảng $0.82$; trên CUB là từ $0.20$--$0.27$ lên $0.85$. Chênh lệch càng lớn ở
        miền khó (fine-grained CUB), nơi học-từ-đầu gần như thất bại.
  \item \textbf{Nguồn cứu đuôi nhất là mô hình nền tảng thị giác thứ hai (DINOv2).} DINOv2-LIFT
        vượt LIFT-CLIP ở cả hai dataset ($0.811$ vs $0.718$; $0.841$ vs $0.670$) và là trụ cột của
        fusion. DINOv2 được huấn luyện tự giám sát với patch $14$, nên bắt chi tiết fine-grained tốt
        hơn CLIP ViT-B/32 ở ảnh nhỏ.
  \item \textbf{Ngôn ngữ (LLM) giúp ít.} CLIP-LLM gần như ngang CLIP zero-shot
        (CIFAR $0.626$ vs $0.627$; CUB $0.484$ vs $0.460$, chỉ nhích nhẹ). Mô tả giàu không bù được
        hạn chế của CLIP ViT-B/32 ở ảnh nhỏ và dữ liệu fine-grained.
  \item \textbf{Sinh dữ liệu và mixup không giúp} (cả hai dataset đều thấp hơn LIFT-CLIP gốc): khi
        đặc trưng đã mạnh, augment thêm chỉ đưa vào nhiễu.
  \item \textbf{Đóng góp của pipeline đo được rõ ràng} (Ablation A): chỉ riêng Balanced Softmax đã
        cứu đuôi từ $\sim0.02$ lên $0.33$ (CIFAR) và từ $0.08$ lên $0.68$ (CUB); thêm adapter cùng
        cosine init đẩy tiếp lên $0.74$/$0.86$.
  \item \textbf{Fine-tune nặng làm hại đuôi} (Ablation B): trên CIFAR, full-FT cho few-shot $0.111$
        so với LIFT đóng băng $0.635$; trên CUB, full-FT còn cho bal-acc thấp nhất ($0.411$) --- một
        biện minh trực tiếp cho thiết kế đóng băng trên cả hai bộ.
  \item \textbf{Bù trừ có nhưng ở mức nhẹ và an toàn.} Fusion (greedy) vượt expert đơn tốt
        nhất ở cả hai dataset ($0.823>0.811$; $0.849>0.841$) mà không tụt ở nhóm nào, vì greedy bảo
        đảm $\ge$ best-single trên validation cho từng nhóm.
\end{enumerate}

\subsection{Điểm mạnh và hạn chế}
Về điểm mạnh, nghiên cứu có một khung thống nhất, đa-dataset, với giao thức chọn-mô-hình-trên-val
nghiêm ngặt; các kết luận nhất quán trên hai miền và hai quy mô lớp; có ablation tách bạch đóng góp
và trung thực với cả các kết quả âm.

Để bảo đảm minh bạch, nhóm nêu rõ một ghi chú thực nghiệm liên quan đến \emph{breakdown} (không ảnh
hưởng tới balanced accuracy hay các kết luận): nhóm from-scratch trên CUB (Bảng~\ref{tab:scratch})
được chấm khi ngưỡng nhóm-shot còn ở mặc định ($100/20$) --- với CUB chỉ có tối đa $50$ ảnh/lớp,
nhóm \emph{many} trở nên rỗng nên cách chia many/medium/few không hợp lệ. Vì balanced accuracy không
phụ thuộc cách chia nhóm nên vẫn dùng để so sánh được, còn cột few của CUB từ-đầu chỉ mang tính tham
khảo; track nền tảng (Bảng~\ref{tab:foundation}) đã dùng đúng ngưỡng $20/10$. Một điểm tích cực về
tính nhất quán: bảng tái dùng ở Track 2 (Bảng~\ref{tab:reuse}) được tính trên đúng các checkpoint
từ-đầu (không phải bản tiền huấn luyện), nên các con số nằm khớp giữa Track 1 và Track 3.
